%% ch12 --- Iterating in the plane: Julia and Mandelbrot
%%
%% Written September 2026. Complex numbers from nothing (a point of the
%% plane, with lengths that multiply and angles that add), the rule
%% z -> z^2 + c and its escape test, the Mandelbrot set, Julia sets, and the
%% payoff that closes Part III: the logistic map is the same rule in
%% disguise, so the real axis of the Mandelbrot set IS Chapter 8's cascade.
%% Every number the reader cannot do in their head comes from
%% code/measure/ch12.py, which imports the listings under code/ch12/ and
%% chaoslab's escape_time and r_to_c.
%%
%% WHAT THIS CHAPTER MAY NOT SAY. It does not derive the period doublings
%% or Feigenbaum's number (Chapter 8 does; here they are found again on the
%% real axis and cited). It does not prove the escape radius or the theorem
%% that the orbit of 0 decides whether a Julia set is in one piece; both are
%% stated and labelled as not proved here. It hands nothing on: it closes
%% Part III.
%%
%% Twin: chapters/pl/ch12-mandelbrot.tex. Section for section, macro for
%% macro, number for number; tools/parity.py compares them.

\chapter{Iterating in the plane: Julia and Mandelbrot}
\label{ch:mandelbrot}
\index{Mandelbrot set}\index{complex numbers}

\noindent
Take a number, square it and add $1$. Do the same to the answer, and again.
From $0$ you get $0$, $1$, $2$, $5$, $26$, and the numbers are gone, off
towards infinity. Add $-1$ instead of $1$ and the same recipe goes $0$,
$-1$, $0$, $-1$ for ever. So here is a question with a yes-or-no answer for
every number $c$ you might add: \emph{started at $0$, does the recipe run
away?}

On the number line the answer is dull: the $c$ that do not run away form
one unbroken stretch, from $-2$ to $\tfrac14$. This chapter asks the same
question of the points of a flat plane, with nothing else changed, and the
answer is the most famous picture in mathematics. It is also
Chapter~\ref{ch:bifurcations}'s road to chaos seen from the side, which is
the last thing Part~III has to show you about the shapes chaos makes.

\begin{outcomes}
\outcome{add and multiply two points of the plane, by hand and with Python's
  complex numbers}
\outcome{run the rule $z \to z^{2} + c$ from a starting point and find the
  step at which its orbit escapes}
\outcome{say what a Julia set and the Mandelbrot set each hold fixed and
  each let vary}
\outcome{convert between the logistic map's $r$ and the Mandelbrot set's
  $c$, and find Chapter~\ref{ch:bifurcations}'s cascade on the real axis}
\outcome{tell the parts of the Mandelbrot set where the orbit settles from
  the parts where it never does}
\end{outcomes}

\section{Numbers that are points}
\label{sec:ch12-plane}
\index{complex numbers!multiplying}

\noindent
Every number in this book so far has lived on a line. A point of a flat
plane needs two numbers: how far across, and how far up. Write the point
three across and four up as $3 + 4\iu$. For now the $\iu$ is only a label
meaning \enquote{up}: the number stuck to it is the up part, and the other
is the across part. Numbers written like this are called \emph{complex
numbers}; the across part is called \emph{real} and the up part
\emph{imaginary}, names you will have reason to distrust by the end of this
section.

Adding two of them is sliding. Add the across parts and the up parts
separately, $(3 + 4\iu) + (1 + 2\iu) = 4 + 6\iu$, exactly as two moves on a
map add up to one.

Multiplying needs a rule, and here it is in words. Every point has a
\emph{length}, its distance from $0$, and an \emph{angle}, how far it is
turned anticlockwise from the positive half of the across line. For
$3 + 4\iu$ the length is $5$, by Pythagoras, and the angle is
\val{ch12.angle} degrees. The rule: \textbf{to multiply two points,
multiply their lengths and add their angles}. Multiplying by $2$, which has
length $2$ and angle $0$, doubles a point's distance from $0$ and leaves its
direction alone. Multiplying by $\iu$, the point one step up, keeps the
length and turns the point a quarter turn anticlockwise.

\begin{think}
The point $\iu$ has length $1$ and angle $90$ degrees. What is
$\iu \times \iu$?
\end{think}
\reveal{$-1$. The lengths multiply to $1 \times 1 = 1$ and the angles add to
$180$ degrees: half a turn from $1$ lands on $-1$.}
School tells you that no number squares to a negative, and on the line that
is true. Off the line it is not, and nothing mysterious has happened: two
quarter turns make a half turn. With $\iu \times \iu = -1$ in hand you can
multiply out brackets as you always have and tidy up at the end:
\[ (a + b\iu)(c + d\iu) = (ac - bd) + (ad + bc)\iu . \]
In words: the across part of the product is the two across parts
multiplied, minus the two up parts multiplied; the up part is each across
part times the other's up part, added. That is the rule a computer uses,
and it agrees with the lengths-and-angles rule for every pair of points.
Python has it built in: it calls the type \code{complex} and writes $\iu$ as
\code{1j}.

\pyregion{code/ch12/complex_numbers.py}{times}{The multiplication rule, by hand}{lst:ch12-times}

\pyregion{code/ch12/complex_numbers.py}{check}{By hand, by Python, and by lengths and angles}{lst:ch12-check}

\transcript{ch12-complex}

\noindent
Square $3 + 4\iu$ both ways. The algebra gives $9 - 16 = -7$ across and
$12 + 12 = 24$ up. The picture gives a length of $5 \times 5 = 25$ and an
angle of \val{ch12.angle.twice} degrees, twice \val{ch12.angle}, and the
point $-7 + 24\iu$ has exactly that length and that angle. Two descriptions,
one answer.

\begin{think}
Numbers with an $\iu$ in them are called imaginary. Does anything in the
world actually behave like $\iu$?
\end{think}
\reveal{Yes: anything that has a size and a direction and can be turned.
Multiplying by $\iu$ is a quarter turn, and turning is as real as anything
you can measure.}

\begin{trapbox}
\textbf{\enquote{Complex numbers are imaginary, so they describe nothing
real.}} The name is history, not a description. A complex number is a point
of the plane with a rule for multiplying, and the rule says how to stretch
and how to turn. Anything that stretches and turns \dash{} a wave, a
rotation, an alternating current \dash{} is described by complex numbers
because that is what they are made for, not because anyone is pretending.
\end{trapbox}

\begin{worldbox}
Electrical engineers describe an alternating voltage as a point going round
a circle: its length is the size of the voltage and its angle is where in
its cycle it is. A circuit multiplies that point by a fixed complex number,
its \emph{impedance}, which scales the size and shifts the timing \dash{}
lengths multiplying, angles adding. Engineers write $j$ for $\iu$, because
$i$ was already taken for current, and Python follows them: that is why it
writes \code{1j}.
\end{worldbox}

\section{Square, add, and ask whether it escapes}
\label{sec:ch12-escape}
\index{escape time}

\noindent
Here is the chapter's whole rule. Pick a complex number $c$. Start at
$z = 0$ and step by
\[ z \to z^{2} + c , \]
where the arrow means \emph{becomes}: square the point, then slide it by
$c$. It is the iteration of Chapter~\ref{ch:iteration}, with a plane in
place of a line. When $c$ sits on the across line this is the opening's
arithmetic: a run to infinity for $c = 1$, a cycle of two for $c = -1$.

\begin{think}
Start at $0$ with $c = -2$ and write down the first few steps. Does the
orbit run away?
\end{think}
\reveal{No. It goes $0$, $-2$, $2$, and then $2 \times 2 - 2 = 2$ for ever:
it sits exactly on $2$ and never gets further out.}

\pyregion{code/ch12/orbits.py}{orbit}{The orbit of 0 under the rule}{lst:ch12-orbit}

\transcript{ch12-orbits}

\noindent
The last row is the plane at work. With $c = \iu$ the orbit leaves the line
at once, runs round a small loop and then swaps between $-1 + \iu$ and
$-\iu$ for ever. The row for $\tfrac14$ creeps upwards and never gets past
$\tfrac12$.

You cannot watch for ever, so you need a test that settles the question
in finitely many steps, and there is one: \textbf{once the
orbit's length passes $2$, it never comes back}. The number of steps it
takes to pass $2$ is the point's \emph{escape time}, and the book's library
computes it in five lines.

\begin{rigourbox}
Why $2$ is enough, in words rather than as a proof. If $c$ is longer than
$2$, the first step already passes $2$, so suppose it is not, and that $z$
is longer than $2$. Squaring makes the length of $z$ its own square, and
adding $c$ pulls a point back by at most the length of $c$, which is shorter
than $z$. So the new length is at least the old length times (the old
length minus one), which is more than the old length: the point moves out,
and the further out it is, the faster it goes. \enquote{Adding $c$ pulls
back by at most its length} is the triangle inequality, proved in any first
course on complex numbers.
\end{rigourbox}

\pyregion{code/src/chaoslab/fractals.py}{escape}{Escape time, from the book's library}{lst:ch12-escape}

\noindent
The test adds the squares of the two parts and compares them with $4$,
which is Pythagoras without the square root: a length passes $2$ exactly
when its square passes $4$. For $c = 1$ the answer is $3$, the step at which
the orbit reaches $5$. A $c$ whose orbit never escapes runs until
\code{limit} and returns it.

\begin{think}
$c = \tfrac14$ never escapes: it is the creeping row of the table. The point
$c = \num{0.26}$ is a hundredth further right. Roughly how many steps does
its orbit take to pass $2$: about $5$, about $30$, or about $300$?
\end{think}
\reveal{About $30$: it passes $2$ at step \val{ch12.slow.260}. It creeps up
towards $\tfrac12$ like its neighbour, lingers there, and only then slips
past.}

\pyregion{code/ch12/slow_escape.py}{slow}{Points close to the edge linger}{lst:ch12-slow}

\transcript{ch12-slow-escape}

\noindent
A little further out an orbit leaves quickly: $c = \num{0.3}$ is gone by
step \val{ch12.slow.300}. But each time $c$ comes ten times closer to
$\tfrac14$, the orbit lingers about
three times as long: \val{ch12.slow.260}, \val{ch12.slow.251},
\val{ch12.slow.2501} steps. So there is no step count after which you can
be sure: a point that has not escaped after a thousand steps may still
escape after ten thousand. Every picture that follows is drawn with a
limit, and close to the edge it is an honest approximation, not a verdict.

\begin{lab}{l12_01_escape}{Escape time without complex numbers}
Store a point as a pair of ordinary numbers and write one step of the rule
with the multiplication rule of Section~\ref{sec:ch12-plane}; then the
escape time on top of it; then two yes-or-no questions: for this $c$, does
the orbit of $0$ stay, and for this fixed $c$, does the orbit of this start
stay? The test checks your escape times against the library's at 480
points of the plane.
\end{lab}

\section{The Mandelbrot set}
\label{sec:ch12-mset}
\index{Mandelbrot set!drawing}

\noindent
Now colour every $c$ in the plane: dark if the orbit of $0$ never escapes,
lighter the sooner it does. The dark points are the \emph{Mandelbrot set}.
The program below draws it in characters, one for each $c$, and counts a
point as in the set if it is still inside after \val{ch12.draw.limit}
steps.

\pyregion{code/ch12/ascii_mandelbrot.py}{draw}{The Mandelbrot set in characters}{lst:ch12-draw}

\noindent\begin{minipage}{\linewidth}
\transcript{ch12-ascii}
\end{minipage}

\noindent
The middle row is the across line, and it is solid \code{@} from $-2$ to
$\tfrac14$, as the opening promised. Off the line the set is a heart-shaped
body with a round bulb to its left, smaller bulbs on both, and a thin
needle running out to $-2$. Figure~\ref{fig:ch12-mandelbrot} draws it
properly, with two squares on its edge magnified.

\plotfig{ch12-mandelbrot}{The Mandelbrot set: black where the orbit of $0$
never escapes, darker blue the longer it lingers before escaping. Square A
sits on the edge where the heart meets its biggest bulb and is shown at two
magnifications; square B sits on the needle.}{the Mandelbrot set and three zooms}

\noindent
Pierre Fatou and Gaston Julia studied rules like this one around 1918 to
1920, with pen and paper and no pictures at all. Benoit Mandelbrot's
computer pictures, published in 1980, made the set famous, and the rule he
iterated in that paper was the logistic rule itself, written with complex
numbers \dash{} which the last section of this chapter explains.

\begin{think}
Magnify the edge a thousand times, a million times, and there is still new
detail: spirals, hairs, tiny copies. Where is all of that detail stored?
\end{think}
\reveal{Nowhere. Every point is worked out afresh from the rule and the
escape test. The program that drew the characters above is
\val{ch12.draw.lines} lines long.}

\begin{trapbox}
\textbf{\enquote{Infinite detail means infinite information.}} The edge has
detail at every scale you can reach, and all of it follows from one line of
algebra and one question. The detail is \emph{generated}, not stored, which
is Chapter~\ref{ch:logistic}'s lesson again: there, one line produced every
kind of behaviour. To tell someone the whole set exactly you need only tell
them the rule and the question; a short description can have an endless
consequence.
\end{trapbox}

\noindent
Some plain facts about the set are still unknown, such as its area. Count
the grid squares whose centre is still inside after 50 steps and you get
about \val{ch12.area.50}; after 200 steps, \val{ch12.area.200}; after 1000,
\val{ch12.area.1000}. The count keeps falling as slow escapers near the edge
are found out, and nobody has a formula for where it stops. The edge itself
is so crumpled that its dimension, in the sense of
Chapter~\ref{ch:fractals}, is $2$, the most that any shape drawn in a plane
can have: Mitsuhiro Shishikura proved it.

\section{Julia sets: fix \texorpdfstring{$c$}{c}, move the start}
\label{sec:ch12-julia}
\index{Julia set}

\noindent
The same rule answers a second question. Hold $c$ fixed and let the
\emph{start} vary: for each point $z$ of the plane, does the orbit that
begins there run away? The starts that stay fill in the \emph{Julia set} of
that $c$, named after Gaston Julia. (Strictly, the Julia set is the edge of
that region.)

\mermaidfig{ch12-two-questions}{One rule, two questions. Fix $c$ and ask
about every start, and you draw a Julia set; start at $0$ and ask about
every $c$, and you draw the Mandelbrot set.}{one rule, two sets}

\begin{think}
Take $c = 0$, so the rule is only \enquote{square the point}. Which starting
points run away?
\end{think}
\reveal{Those more than $1$ from the centre. Squaring squares the length: a
length above $1$ grows without end, one below $1$ shrinks towards $0$, and a
length of exactly $1$ stays $1$. The starts that stay fill a disc of radius
$1$, and its edge is a circle.}

\noindent
Change $c$ and the circle crumples. Figure~\ref{fig:ch12-julia} shows two. For $c = -1$ the set is in one
piece, a chain of bulbs pinched together. For $c = \num{-0.8} +
\num{0.2}\iu$ it has fallen apart into dust: infinitely many specks with no
two joined, like the Cantor set of Chapter~\ref{ch:fractals}.

\plotfig{ch12-julia}{Two Julia sets, black where a start never escapes. For
$c = -1$ the set is one piece; for $c = \num{-0.8} + \num{0.2}\iu$ only dust
is left, too fine to show as black, marked by the dark spirals where starts
linger longest.}{a Julia set in one piece and one in dust}

\noindent
What decides which? The orbit of $0$, which is the Mandelbrot question. For
$c = -1$ it swaps between $0$ and $-1$ and stays; for $c = \num{-0.8} +
\num{0.2}\iu$ it passes $2$ at step \val{ch12.dust.escape}. A theorem, not
proved here but in any book on complex dynamics, says that this always
decides it: the Julia set is in one piece exactly when $c$ is in the
Mandelbrot set, and is dust exactly when it is not. So the Mandelbrot set is
a map of every Julia set at once, and each black point of it is a Julia set
that holds together.

\section{The logistic map in disguise}
\label{sec:ch12-logistic}
\index{logistic map!and the Mandelbrot set}

\noindent
The chapter's rule and Chapter~\ref{ch:logistic}'s logistic map
$x \to rx(1 - x)$ look different, and they are the same rule. Relabel the
line: measure each $x$ from $\tfrac12$ instead of from $0$, flip it, and
stretch it by $r$, which is the change
\[ z = r\left(\tfrac12 - x\right) . \]
Then $x(1 - x) = \tfrac14 - z^{2}/r^{2}$, and the next $z$, which is
$\tfrac{r}{2} - r^{2}x(1 - x)$, comes out as $z^{2} + c$ with
\[ c = \frac{r}{2} - \frac{r^{2}}{4} . \]
In words: halve $r$, then take away a quarter of its square. Nothing is lost
in the translation. Every orbit of one rule is an orbit of the other on a
relabelled line, and $x = \tfrac12$ becomes $z = 0$: the orbit of $0$ that
the Mandelbrot set asks about is the logistic orbit from $\tfrac12$.

\begin{think}
Chapter~\ref{ch:logistic}'s fixed point loses its grip at $r = 3$, and a
cycle of two takes over. Which $c$ is that?
\end{think}
\reveal{$c = \tfrac32 - \tfrac94 = \num{-0.75}$: exactly where the heart
meets the round bulb to its left.}
So the heart is the fixed point and that bulb is the cycle of two. As $r$
runs from $1$ to $3$, $c$ runs from $\tfrac14$ down to $\num{-0.75}$ along the
heart's stretch of the line. The listing runs both rules side by side at
the values of $r$ from Chapter~\ref{ch:logistic}, from starts that the
change matches up, and reports the length of the cycle each settles into.

\pyregion{code/ch12/same_rule.py}{same}{The same rule, run twice}{lst:ch12-same}

\transcript{ch12-same-rule}

\noindent
The cycles agree at every $r$ \dash{} one, two, four, eight, three, and none
at all at $\num{3.9}$ \dash{} and that is Chapter~\ref{ch:bifurcations}'s
bifurcation diagram written along the real axis.
Figure~\ref{fig:ch12-real-axis} sets one picture above the other.

\plotfig{ch12-real-axis}{Above, the Mandelbrot set along the real axis;
below, Chapter~\ref{ch:bifurcations}'s bifurcation diagram with each $r$
drawn at its $c$. Every splitting of the cycle happens where one bulb meets
the next. The dashed lines mark $r = 3$, $r = 1 + \sqrt{6}$, the Feigenbaum
point and the period-three window.}{the real axis over the bifurcation diagram}

\noindent
The second splitting, at $r = 1 + \sqrt{6}$, converts to
$c = \val{ch12.c.two}$, where the first bulb meets the next. Each bulb also
has a centre, the $c$ at which the orbit of $0$ comes back to exactly $0$:
$-1$, then \val{ch12.centre.4}, then \val{ch12.centre.8}, and on. Converted
to $r$ they are \val{ch12.rc.2}, \val{ch12.rc.4} and \val{ch12.rc.8}: the
values at which Chapter~\ref{ch:bifurcations}'s cycles pass through
$\tfrac12$, its superstable parameters. The gaps between the centres shrink
by a factor that settles at \val{ch12.feig.ratio}, Feigenbaum's number
again, and the bulbs pile up at $c = \val{ch12.feig.c}$, which is
$r = \val{ch12.feig.r}$: the Feigenbaum point. Beyond it lies the needle.
Even the period-three window is there. It opens at $c = \val{ch12.c.window}$,
and square B of Figure~\ref{fig:ch12-mandelbrot} shows what sits on the
needle at that point: a tiny copy of the whole set.

\begin{think}
Chapter~\ref{ch:logistic}'s orbits that never settle live at $r =
\num{3.9}$. The table above converts it. Is that $c$ in the heart, in a
bulb, or on the needle?
\end{think}
\reveal{On the needle: $c$ is about $\num{-1.85}$, left of the Feigenbaum
point, past the last of the bulbs that pile up there.}

\begin{trapbox}
\textbf{\enquote{The Mandelbrot set is a picture of chaos.}} Most of what is
black is order. In the heart the orbit of $0$ settles on one point, in the
big bulb on a cycle of two, and in every bulb on a cycle whose length the
bulb fixes. The never-settling orbits of Chapter~\ref{ch:logistic} sit on
the thin edge, and on the real axis that means the needle, between its
windows. The famous picture is mostly a map of order, with chaos living
along its border.
\end{trapbox}

\begin{lab}{l12_02_convert}{From $r$ to $c$ and back}
Write the conversion both ways, $r$ to $c$ and $c$ back to $r$, and the
change $z = r(\tfrac12 - x)$ itself. The test runs the logistic map and
$z \to z^{2} + c$ side by side and checks that your change keeps them in
step. Then send the period-doubling values of $r$ in the starter across and
see where on the real axis they land.
\end{lab}

\noindent
That closes Part~III: chaos has a shape, and its last picture contains
Chapter~\ref{ch:bifurcations} whole. Part~IV takes the ideas out into
pendulums, planets, weather and living things.

\begin{summarybox}
\begin{itemize}
\item A \textbf{complex number} is a point of the plane, written $a +
  b\iu$. Adding is sliding; multiplying multiplies the \textbf{lengths} and
  adds the \textbf{angles}, the same as multiplying out with
  $\iu \times \iu = -1$.
\item The rule $z \to z^{2} + c$, run from $0$, either stays or escapes, and
  once the length passes $2$ it has escaped for good. The \textbf{escape
  time} counts the steps; near the edge it grows without limit.
\item The \textbf{Mandelbrot set} is every $c$ whose orbit of $0$ never
  escapes: a heart, bulbs, a needle to $-2$, and an edge with detail at
  every scale, generated by the rule. Its area is not known exactly.
\item A \textbf{Julia set} fixes $c$ and asks about every start. It is in
  one piece when $c$ is in the Mandelbrot set and dust when it is not.
\item The change $z = r(\tfrac12 - x)$ turns the \textbf{logistic map} into
  $z \to z^{2} + c$ with $c = \tfrac{r}{2} - \tfrac{r^{2}}{4}$, so the real
  axis of the set is the bifurcation diagram: bulbs are cycles, they pile up
  at the Feigenbaum point, and chaos lives on the needle.
\end{itemize}
\end{summarybox}

\canyou

\begin{checkpoint}
\item Work out $(2 + 3\iu) + (1 - \iu)$ and $(1 + \iu)(1 - \iu)$.
  \answerto{$3 + 2\iu$, and $2$: multiplying out gives
  $1 - \iu + \iu - \iu \times \iu = 1 + 1$.}
\item Use lengths and angles to find $(1 + \iu)^{2}$, then check it by
  multiplying out.
  \answerto{$1 + \iu$ has angle $45$ degrees, so its square has angle $90$
  and length $2$ (the length of $1 + \iu$, squared): it is $2\iu$. Multiplying
  out, $1 + 2\iu + \iu \times \iu = 2\iu$.}
\item Start at $0$ with $c = 2$. At which step does the orbit pass $2$?
  \answerto{At step $2$: the orbit is $0$, $2$, $6$, and $2$ is not past $2$
  but $6$ is.}
\item Is $c = \num{-1.5}$ in the Mandelbrot set?
  \answerto{Yes. Every real $c$ from $-2$ to $\tfrac14$ is in it; the orbit
  of $0$ stays between $-2$ and $2$.}
\item For $c = 0$, does the start $\num{0.6} + \num{0.9}\iu$ run away?
  \answerto{Yes. Its length is the square root of $\num{0.36} + \num{0.81}$,
  a little more than $1$, and squaring keeps making a length above $1$
  larger.}
\item Convert $r = \num{3.2}$ to $c$. Is it in the heart or in the big bulb,
  and does that agree with Chapter~\ref{ch:logistic}?
  \answerto{$c = \num{1.6} - \num{2.56} = \num{-0.96}$, inside the big bulb,
  which runs from $\num{-0.75}$ to $\num{-1.25}$: a cycle of two, as
  Chapter~\ref{ch:logistic} found at that $r$.}
\item Why is the Mandelbrot set not a picture of chaos?
  \answerto{Its heart and bulbs are values of $c$ at which the orbit of $0$
  settles into a cycle. The chaotic values sit on the thin edge, which on
  the real axis is the needle between its windows.}
\end{checkpoint}
